% TOP_PROBLEM: {"id": 8400004, "title": "O5a — Deterministic polynomial-time integer factorization", "category_id": 3, "category": {"id": 3, "name": "graph_theory", "display_name": "Graph Theory", "description": "Problems involving graphs, networks, and their properties.", "slug": "graph-theory", "order_index": 3, "created_at": "2026-07-31T15:26:25.671Z"}, "status": "open"}
% FIRST_POSTED: 2026-09-27
% !TeX program = lualatex
% ATTEMPT_STATUS: unresolved
% Research continuation by GPT_6_Astra_Ultra, 27 September 2026.
% Catalogue statements and existing sources preserved verbatim.
\documentclass[11pt]{article}
\usepackage[margin=25mm]{geometry}
\usepackage{fontspec}
\setmainfont{Latin Modern Roman}
\usepackage{amsmath,amssymb,mathtools}
\usepackage{unicode-math}
\setmathfont{Latin Modern Math}
\usepackage{xurl,hyperref}
\hypersetup{colorlinks=true,urlcolor=blue}
\setcounter{secnumdepth}{0}
\setlength{\parindent}{0pt}
\setlength{\parskip}{5pt}
\setlength{\emergencystretch}{3em}
\begin{document}
\section*{\texorpdfstring{O5a — Deterministic polynomial-time integer factorization}{O5a — Deterministic polynomial-time integer factorization}}\label{8400004}
\subsection{Definitions and mathematical statement}
The binary length of $N\ge2$ is $\ell(N)=\lfloor\log_2N\rfloor+1$. A nontrivial factor is an integer $d$ with $1<d<N$ and $d\mid N$. Ask whether there exist a deterministic classical algorithm $A$ and constants $C,k$ such that
\[
 \forall N\text{ composite},\qquad A(N)=d\text{ a nontrivial factor},\qquad
 T_A(N)\le C\ell(N)^k.
\]
The promise is compositeness; the running-time bound is in the bit length, not in the value of $N$.
\par\smallskip\textit{Formulation and concept references:} \href{https://www.sciencedirect.com/science/article/pii/S0022000015000768}{[S1]}.

\subsection{Short English statement}
Is there a deterministic classical method that factors every composite integer in time polynomial in its number of binary digits?
\subsection{Research attempt}

\paragraph{Scope and checked algorithmic context.}
The input is promised composite, and the requested output is one proper
divisor. Running time means bit operations in $\ell=\ell(N)$; a loop of
length $N^{1/5}$ remains exponential in this parameter. An established
deterministic result of Harvey--Hittmeir [E1] factors $N$ in
$O(N^{1/5}\log^{16/5}N/(\log\log N)^{3/5})$ bit operations. It is
substantially faster than trial division, but does not meet this target.
Their 2026 manuscript [E2] improves the auxiliary problem of finding
large-order elements; its stated alternatives include returning an element,
so it is not a polynomial-time factorization theorem. No completeness
claim about all current algorithms is needed here.

The chosen route examines modular square roots. Je\v r\'abek [E3], the
author manuscript of [S1], studies reductions involving factoring and
square-root search, carefully distinguishing randomized reductions from
deterministic ones. We give the elementary fiber calculation behind this
connection, then test the tempting derandomization by a short fixed list
of inputs. The calculation is classical; its purpose here is to expose
the exact unsupplied algorithmic step rather than rebrand a reduction as
a factorization algorithm. The conditional claims in catalogue source
[S2] are not assumed.

\paragraph{Elementary preprocessing is genuinely polynomial.}
If $N$ is even, return $2$. Detect perfect powers by trying each integer
$2\le k\le\ell$ and computing the integer $k$th root by binary search,
with exact exponentiation and comparison to $N$. Intermediate products
may be capped as soon as they exceed $N$, so the bit lengths remain
$O(\ell)$. There are $O(\ell)$ exponents, $O(\ell)$ search steps per
exponent, and polynomially many bit operations per step. If $N=a^k$,
$a>1$, return $a$. This includes odd prime powers without needing to
recognize their prime base. The unhandled promised inputs are odd and
have at least two distinct prime factors.

Write their unknown factorization, for analysis only, as
$N=\prod_{i=1}^r p_i^{e_i}$, with $r\ge2$. This notation does not give
the algorithm access to the factors. Let $U_N=(\mathbb Z/N\mathbb Z)^\times$
and let $Q_N$ be the image of its squaring map. A hypothetical oracle
$F_N$ takes $u\in Q_N$ and returns any $v$ satisfying $v^2\equiv u\pmod N$.
Because $u$ is a unit, its returned root is automatically a unit. The
oracle is not presumed to return a random root or a preferred CRT sign.

\paragraph{Exact root fibers and extraction probability.}
\textbf{Lemma 17.1.} Every $u\in Q_N$ has exactly $2^r$ square roots
modulo $N$. For a fixed root $v$, exactly two roots $a$ fail to give a
proper factor through $\gcd(a-v,N)$.

\textit{Proof.}
For an odd prime power $p^e$, a root $w$ of $w^2=1$ satisfies
$p^e\mid(w-1)(w+1)$. Since $\gcd(w-1,w+1)$ divides two, the full
prime power divides precisely one factor. Thus $w\equiv1$ or $-1$
modulo $p^e$. Multiplication by any fixed root of $u$ bijects its roots
with roots of one. The Chinese remainder theorem therefore gives $2^r$
roots, one for each independent local sign.

If $a$ uses the same sign as $v$ at every prime power, then $a=v$
modulo $N$ and the gcd is $N$. If every sign is opposite, then $a=-v$
and the gcd is one, since $2v$ is a unit. Any mixed sign pattern makes
the gcd the product of a nonempty proper subset of the prime-power
factors. This is a proper divisor, and there are exactly $2^r-2$ such
patterns.\hfill$\square$

Choose $a$ uniformly in $U_N$ and query $F_N(a^2)$. Conditional on the
square $u=a^2$, the element $a$ is uniform in its $2^r$-element fiber.
The oracle's chosen root depends on $u$, not on which preimage was
sampled. Hence the success probability is exactly
\begin{equation}\label{ra:17:probability}
 1-2^{1-r}\ge\tfrac12.
\end{equation}
This remains true for any fixed deterministic selection of one root
per square. If the oracle itself is randomized independently of the
sampled preimage, condition on its output and obtain the same result.

An implementable sampler need not already know $U_N$. Sample uniformly
from $1,\ldots,N-1$ and compute $g=\gcd(a,N)$. If $g>1$, it is already
proper and succeeds. Otherwise the conditional distribution is uniform
on $U_N$ and the lemma applies. Uniform bounded-integer sampling by
rejection uses expected $O(\ell)$ random bits. Thus a polynomial-time
square-root oracle would yield a randomized expected-polynomial-time
factorization algorithm, with failure probability at most $2^{-t}$
after $t$ independent trials. Every reported factor is checked by exact
division. This is not the deterministic worst-case bound requested.

\paragraph{Attempted derandomization and an exact obstruction.}
A natural proposal is to replace the random $a$ by the first
$\operatorname{poly}(\ell)$ integers, or another short deterministic
list $H\subseteq U_N$. A success proportion of one half on the whole
unit group does not imply that this particular list meets the successful
half. Here the possible obstruction can be characterized exactly.

\textbf{Lemma 17.2.} There exists a valid root selector $F_N$ for which
every test $a\in H$ fails if and only if, in each squaring fiber, the
elements of $H$ lie in a single pair $\{b,-b\}$.

\textit{Proof.}
If that condition holds, for each fiber meeting $H$ select one member
$b$ of its represented sign pair as the oracle output. Then all tested
preimages equal $b$ or $-b$, and Lemma 17.1 says they fail. Choose any
root on unqueried fibers to complete a valid selector. Conversely, if
a fiber contains tested elements $a,b$ with $a\not\equiv\pm b$, no
selected root can have both of them in its single exceptional sign pair.
At least one query succeeds.\hfill$\square$

In fact, in the latter case the oracle is unnecessary:
$a^2\equiv b^2\pmod N$ and $a\not\equiv\pm b$ already imply that
$\gcd(a-b,N)$ is a proper divisor by the same local-sign argument.
Consequently a short list guaranteed to defeat \emph{every} root selector
must already contain a useful congruence of squares. Producing such lists
deterministically in polynomial time would itself solve the remaining
factoring instances. The random reduction has not made that construction
automatic.

The obstruction also applies to a finite adaptive transcript: when a
previously unseen square is queried, an adversary can return the supplied
preimage; when a previously queried square recurs in the same sign pair,
the old response still causes failure. A new sign pair exposes precisely
the factor-producing collision. This adversary can be extended to a
valid selector after the transcript. Its implementation need not be
efficient, so this is \emph{not} a complexity lower bound against efficient
root algorithms, nor does it exclude exploiting the structure of a
specific algorithm. It isolates the defect in an oracle-independent
derandomization based only on the success fraction.

\paragraph{A concrete failed list and a successful collision.}
For $N=77$, the squares of $1,2,3,4,5,6$ are distinct modulo $77$,
and all six inputs are units. A valid selector may return the respective
input on those six squares; all six gcd tests then fail. Nevertheless
there are four roots of one,
$1,34,43,76$, and the mixed roots immediately yield
$\gcd(34-1,77)=11$ and $\gcd(34+1,77)=7$.
The roots exist in abundance relative to the two trivial choices, but
their existence does not tell the deterministic procedure where to find
one. This small example illustrates the logical failure of the proposed
list argument, not an asymptotic hard instance: trial division already
factors $77$ quickly.

One can also defeat a proposed use of generic collision finding. The
squaring map has $\varphi(N)/2^r$ images, and evaluating it is cheap.
Enumerating enough preimages to force a collision by the pigeonhole
principle can still require exponentially many in $\ell$. Moreover the
automatic collision $a,-a$ gives no factor. A collision theorem must
control both its search cost and its local sign pattern. Calling a
collision algorithm polynomial in $N$, or counting arithmetic operations
while allowing exponentially long integers, would not meet the stated
bit-complexity target.

\paragraph{Verification and edge cases.}
The local script enumerates every unit and every squaring fiber for
$N=15,45,77,105,315,1001$. For the selector choosing the smallest root,
the exact success counts are
\[
\begin{array}{r|r|r|r}
N&r&\varphi(N)&\text{successful units}\\\hline
15&2&8&4\\45&2&24&12\\77&2&60&30\\
105&3&48&36\\315&3&144&108\\1001&3&720&540
\end{array}
\]
The repeated-prime cases $45$ and $315$ check that the lemma uses distinct
prime factors, not the total number counted with multiplicity. For every
unit it verifies both the gcd criterion and the exact fraction in
\eqref{ra:17:probability}. The table-building oracle is implemented by
exhaustive enumeration and has no claimed efficient running time.
The script separately verifies the six-query failure setup for $77$.
These are exact integer calculations, stored in \texttt{checks.py} and
\texttt{results.json} under
\texttt{\_Local/top\_problems/continuation\_15\_25\_2026\_09\_27/analytic/}.

Even moduli were removed before the odd-prime-power argument. A pure
prime power would have $r=1$ and success probability zero, which is why
the perfect-power step is essential. On a prime input no proper divisor
exists; the catalogue promises compositeness, and no total running-time
claim on off-promise inputs is silently used. On the actual promised
inputs, gcds, modular squaring, and output verification are polynomial;
the missing operations are finding roots and removing randomness.

\paragraph{Remaining gap and outcome.}
The fully proved result is the exact fiber reduction and its sharp
failure condition for a deterministic test set. A complete solution
along this route needs an actual deterministic polynomial-time root
procedure together with a polynomial-time way to force a mixed-sign
collision, or a different direct divisor construction. Neither is
provided. An impossibility claim would need a lower bound against all
deterministic algorithms, far beyond the limited oracle observation.

\textbf{Outcome: UNRESOLVED.} The randomized reduction and the oracle
stress test are justified; no deterministic polynomial-time factoring
algorithm or superpolynomial lower bound has been obtained. The process
combined direct proof, primary-source checking, and exact finite tests;
the results are not a claim of historical novelty or external certification.

\subsection{Sources}
\begin{itemize}
\item[S1]Integer factoring and modular square roots.
\url{https://www.sciencedirect.com/science/article/pii/S0022000015000768}.
\textit{Formulation source.}
\item[S2]A number-theoretic conjecture implying faster algorithms for polynomial factorization and integer factorization — Chris Umans and Siki Wang (2025).
\url{https://arxiv.org/abs/2511.10851}.
\textit{Further reference; not a proof endorsement.}
\item[S3]Improved Separation Between Quantum and Classical Computers for Sampling and Functional Tasks — Marshall, Aaronson and Dunjko, CCC 2025.
\url{https://drops.dagstuhl.de/storage/00lipics/lipics-vol339-ccc2025/LIPIcs.CCC.2025.5/LIPIcs.CCC.2025.5.pdf}.
\textit{Further reference; not a proof endorsement.}

\item[E1]David Harvey and Markus Hittmeir,
\emph{A log-log speedup for exponent one-fifth deterministic integer
factorisation}, Mathematics of Computation 91 (2022), 1367--1379;
arXiv:2105.11105.
\url{https://arxiv.org/abs/2105.11105}.
\textit{Primary research reference for a proved deterministic bit-complexity
bound, which is exponential in input length. Consulted 27 September 2026.}
\item[E2]David Harvey and Markus Hittmeir,
\emph{Deterministic methods for finding elements of large multiplicative order},
arXiv:2601.11131v2 (5 June 2026).
\url{https://arxiv.org/abs/2601.11131v2}.
\textit{Recent manuscript: auxiliary order-search result, not a claimed
polynomial-time factoring algorithm; not used in the proofs here.
Consulted 27 September 2026.}
\item[E3]Emil Je\v r\'abek, \emph{Integer factoring and modular square roots},
arXiv:1207.5220v3 (28 July 2015), author manuscript of [S1].
\url{https://arxiv.org/pdf/1207.5220v3}.
\textit{Primary source on factoring and square-root search reductions;
the elementary fiber reduction is proved explicitly above.
Consulted 27 September 2026.}
\end{itemize}
\end{document}
